\documentclass[12pt, letterpaper, fleqn, leqno]{article}
\usepackage{graphicx} % Required for inserting images
%\usepackage[T1]{fontenc}

\usepackage{polski}
%\usepackage[utf8]{inputenc}

\usepackage{hyperref}
\usepackage{amsmath}
\usepackage[inline]{enumitem}
\usepackage{setspace}
\onehalfspacing

\title{Programy użytkowe \\ Lista 2}
\author{ }
\date{}

\begin{document}
\maketitle

\begin{enumerate}
    % \item Listy tworzymy przy użyciu polecenia
    % \begin{verbatim}
    %     \begin{enumerate}
    %         \item ...
    %     \end{enumerate}
    % \end{verbatim}
    % Wszystko, co następuje po \verb|\item| należy do danego punktu. Rozwiązania tej listy umieść w kolejnych punktach (jeśli ma to zastosowanie). 
    % \item
    \item Wprowadź równanie w trybie wyróżnionym, a następnie nadaj mu etykietę (\verb|\label|) oraz użyj odniesienia do tego równania(\verb|\ref|, lub przy użyciu pakietu \verb|amsmath| polecenia \verb|\eqref|) w tekście przed lub po równaniu. Załaduj pakiet \verb|hyperref| i sprawdź efekt.
    \item Nie tylko do równań możemy tworzyć odnośniki. Wstaw do dokumentu obrazek i nadaj mu etykietę. Odwołaj się do niego w tekście poniżej.
    \item Równania wielolinijkowe możemy (między innymi) tworzyć przy pomocy poleceń \verb|\split|
    \begin{verbatim}
    \begin{equation}
            \begin{split}
            \end{split}
    \end{equation}
    \end{verbatim}
    oraz \verb|\align|. Wypróbuj każde z nich do wpisania poniższego wyrażenia
    \begin{align*}
       \sum_{i=1}^n i   = & 1+2+\dots+n\\
        = & \frac{n(n+1)}{2}.
    \end{align*}
    i zaobserwuj różnicę w ich działaniu. Zwróć uwagę na konieczność załadowania odpowiedniego pakietu.
    \item Przy pomocy  \verb|cases| zapisz poniższe wyrażenia
    
    \begin{enumerate}
        \item \begin{equation}
            \begin{cases}
                x+y=3\\
                x-2y=0
            \end{cases}
        \end{equation}
        \item \begin{equation} 
        \sqrt{x^2} = \begin{cases}
                x, & \text{gdy } x\geq 0,\\
                -x, & \text{gdy} x<0.
            \end{cases}
        \end{equation}
        Zwróć uwagę, na wyrównanie tekstu opisującego warunek (wskazówka: użyj \verb|&|), oraz na czcionkę (\verb|\text{}| wewnątrz trybu matematycznego.).
    \end{enumerate}
    \item Wprowadź w dokumencie kilka akapitów Lorem Ipsum (\verb|\lipsum|). Poeksperymentuj z układem tekstu: marginesami (pakiet \verb|geometry|), interlinią (pakiet \verb|setspace|) oraz czcionką i układem kolumn (opcje w klasie \verb|article|).
\end{enumerate} 
%
\textbf{Zadanie domowe.} Zapisz w \LaTeX ie wybrane polecenie z listy zadań z dowolnego przedmiotu matematycznego, które zawiera w treści jakąś formułę matematyczną, najlepiej rozwiązane przez Ciebie przy tablicy (jeżeli zadanie ma wiele podpunktów, możesz wybrać jeden).

\newpage
\begin{equation}
    \frac{\mathrm{d}x}{\mathrm{d}y} = \frac{\operatorname{d}x}{\operatorname{d}y}
\end{equation}
\begin{equation}
    \operatorname{arg\;max} f(x)
\end{equation}
\end{document}